\documentclass{../templateReading/cdreading}

\newcommand{\authorinfo}{THE FORUM CENTER - NGUYỄN HOÀNG HUY}
\newcommand{\documenttitle}{TÀI LIỆU CHUYÊN ĐỀ READING}
\newcommand{\collectiontitle}{Level 2 — W5 (Jan 26-Feb 1)}

\begin{document}

\thispagestyle{firstpage}
\makeworkshopheader{\authorinfo}{\documenttitle}{\collectiontitle}
\pagestyle{otherpages}

\cdsection{Decision making and Happiness}

\textbf{A}\par

Americans today choose among more options in more parts of life than has ever been possible before. To an extent, the opportunity to choose enhances our lives. It is only logical to think that if some choice is good, more is better; people who care about having infinite options will benefit from them, and those who do not can always just ignore the 273 versions of cereal they have never tried. Yet recent research strongly suggests that, psychologically, this assumption is wrong. Although some choice is undoubtedly better than none, more is not always better than less.

\textbf{B}\par

Recent research offers insight into why many people end up unhappy rather than pleased when their options expand. We began by making a distinction between “maximizers” (those who always aim to make the best possible choice) and “satisficers” (those who aim for “good enough,” whether or not better selections might be out there).

\textbf{C}\par

In particular, we composed a set of statements-the Maximization Scale-to diagnose people’s propensity to maximize. Then we had several thousand people rate themselves from 1 to 7 (from “completely disagree” to “completely agree”) on such statements as “I never settle for second best.” We also evaluated their sense of satisfaction with their decisions. We did not define a sharp cutoff to separate maximizers from satisficers, but in general, we think of individuals whose average scores are higher than 4 (the scale’s midpoint) as maximizers and those whose scores are lower than the midpoint as satisficers. People who score highest on the test – the greatest maximisers-engages in more product comparisons than the lowest scorers, both before and after they make purchasing decisions, and they take longer to decide what to buy. When satisficers find an item that meets their standards, they stop looking. But maximizers exert enormous effort to read labels, checking out consumer magazines and trying new products. They also spend more time comparing their purchasing decisions with those of others.

\textbf{D}\par

We found that the greatest maximizers are the least happy with the fruits of their efforts. When they compare themselves with others, they get little pleasure from finding out that they did better and substantial dissatisfaction from finding out that they did worse. They are more prone to experiencing regret after purchase, and if their acquisition disappoints them, their sense of well-being takes longer to recover. They also tend to brood or ruminate more than satisficers do.

\textbf{E}\par

Does it follow that maximizers are less happy in general than satisficers? We tested this by having people fill out a variety of questionnaires known to be reliable indicators of well-being. As might be expected, individuals with high maximization scores experienced less satisfaction with life and were less happy, less optimistic and more depressed than people with low maximization scores. Indeed, those with extreme maximization ratings had depression scores that placed them in the borderline clinical range.

\textbf{F}\par

Several factors explain why more choice is not always better than less, especially for maximizers. High among these are “opportunity costs.” The quality of any given option cannot be assessed in isolation from its alternatives. One of the “costs” of making a selection is losing the opportunities that a different option would have afforded. Thus an opportunity cost of vacationing on the beach in Cape Cod might be missing the fabulous restaurants in the Napa Valley. Early decision-making research by Daniel Kahneman and Amos Tversky showed that people respond much more strongly to losses than gains. If we assume that opportunity costs reduce the overall desirability of the most preferred choice, then the more alternatives there are, the deeper our sense of loss will be and the less satisfaction we will derive from our ultimate decision.

\textbf{G}\par

The problem of opportunity costs will be worse for a maximizer than for a satisficer. The latter’s “good enough” philosophy can survive thoughts about opportunity costs. In addition, the “good enough” standard leads to much less searching and inspection of alternatives than ‘the maximizer’s “best” standard. With fewer choices under consideration, a person will have fewer opportunity costs to subtract.

\textbf{H}\par

Just as people feel sorrow about the opportunities they have forgone, they may also suffer regret about the option they settle on. My colleagues and I devised a scale to measure proneness to feeling regret, and we found that people with high sensitivity to regret are less happy, less satisfied with life, less optimistic and more depressed than those with low sensitivity. Not surprisingly, we also found that people with high regret sensitivity tend to be maximizers. Indeed, we think that worry over future regret is a major reason that individuals become maximizers. The only way to be sure you will not regret a decision is by making the best possible one. Unfortunately, the more options you have and the more opportunity costs you incur, the more likely you are to experience regret.

I

In a classic demonstration of the power of sunk costs, people were offered season subscriptions to a local theater company. Some were offered the tickets at full price and others at a discount. Then the researchers simply kept track of how often the ticket purchasers actually attended the plays over the course of the season. Full-price payers were more likely to show up at performances than discount payers. The reason for this, the investigators argued, was that the full-price payers would experience more regret if they did not use the tickets because not using the more costly tickets would constitute a bigger loss. To increase the sense of happiness, we can decide to restrict our options when the decision is not crucial. For example, make a rule to visit no more than two stores when shopping for clothing.

\cdsection{Questions 28-31}

Use the information in the passage to match the category (listed A-D) with descriptions or deeds below.

\cdnoteinline{Write the appropriate letters A-D in boxes 28-31 on your answer sheet.}

A     Maximiser

B     Satisficer

C     Both

D     Neither of them

\textbf{28} finish transaction when the items match their expectation

\textbf{29} buy the most expensive things when shopping

\textbf{30} consider repeatedly until they make a final decision

\textbf{31} participate in the questionnaire of the author

\cdsection{Questions 32-36}

Do the following statements agree with the information given in Reading Passage 3?

In boxes 32-36 on your answer sheet, write

TRUE               if the statement is true

FALSE              if the statement is false

NOT GIVEN    if the information is not given in the passage

\textbf{32} With society’s advancement, more chances make our lives better and happier.

\textbf{33} There is a difference of findings by different gender classification.

\textbf{34} The feeling of loss is greater than that of acquisition.

\textbf{35} ‘Good enough’ plays a more significant role in pursuing ‘best’ standards of the maximizer.

\textbf{36} There are certain correlations between the “regret” people and the maximisers.

\cdsection{Questions 37-40}

\cdnoteinline{Choose the correct letter, A, B, C or D.}

\cdnoteinline{Write your answers in boxes 37-40 on your answer sheet.}

\textbf{37} What is the subject of this passage?

A   regret makes people less happy

B   choices and Well-being

C   an interesting phenomenon

D   advices on shopping

\textbf{38} According to the conclusion of questionnaires, which of the following statement is correct?

A   maximisers are less happy

B   state of being optimistic is important

C   uncertain results are found.

D   maximisers tend to cross the bottom line

\textbf{39} The experimental on theater tickets suggested:

A   sales are different according to each season

B   people like to spend on the most expensive items

C   people feel depressed if they spend their vouchers

D   people will feel regret more when they fail to use a higher-priced purchase

\textbf{40} What is the author’s suggestion on how to increase happiness:

A   focus on the final decision

B   be sensitive and smart

C   reduce the choice or option

D   read the label carefully

\cdblue{\textbf{Hướng dẫn giải}}

\cdsection{Phần I: Câu hỏi 28-31 (Matching Features)}

Mục tiêu: Phân loại các hành động hoặc đặc điểm tương ứng với các nhóm đối tượng (A–D).
\begin{itemize}
\item A: Maximiser (Người tối đa hóa)
\item B: Satisficer (Người biết đủ)
\item C: Both (Cả hai)
\item D: Neither of them (Không ai trong số họ)
\end{itemize}

\cdstep{1. Chiến lược tiếp cận}
\begin{itemize}
\item Bước 1: Phân tích câu hỏi để tìm từ khóa chỉ hành động.
\item Bước 2: Định vị thông tin trong bài dựa trên các từ khóa liên quan đến hành vi mua sắm hoặc ra quyết định.
\item Bước 3: Đối chiếu hành vi đó với định nghĩa của "Maximiser" và "Satisficer" trong bài để chọn đáp án.
\end{itemize}

\cdstep{2. Phân tích chi tiết từng câu hỏi}
\begin{itemize}
\item Câu 28: “finish transaction when the items match their expectation”
\begin{itemize}
\item Từ khóa: "finish transaction", "items match expectation".
\item Định vị: Đoạn C, câu: "When satisficers find an item that meets their standards, they stop looking."
\item Đối chiếu: Cụm từ "stop looking" (ngừng tìm kiếm) tương đương với "finish transaction" (hoàn tất giao dịch). Cụm "meets their standards" (đạt tiêu chuẩn) tương đương với "match their expectation" (khớp với mong đợi). Hành động này thuộc về Satisficers.
\item Đáp án: B
\end{itemize}
\item Câu 29: “buy the most expensive things when shopping”
\begin{itemize}
\item Từ khóa: "buy", "most expensive things".
\item Định vị: Bài đọc có nhắc đến việc so sánh sản phẩm (product comparisons), đọc nhãn (read labels) ở đoạn C, hoặc vé xem kịch giá trọn gói (full price) ở đoạn I.
\item Đối chiếu: Tuy nhiên, không có đoạn nào khẳng định rằng Maximiser hay Satisficer có thói quen "luôn mua những thứ đắt nhất". Maximiser chỉ bỏ nhiều công sức tìm kiếm cái tốt nhất (best possible choice), không nhất thiết là cái đắt nhất.
\item Kết luận: Không có thông tin gán hành động này cho nhóm nào cụ thể.
\item Đáp án: D
\end{itemize}
\item Câu 30: “consider repeatedly until they make a final decision”
\begin{itemize}
\item Từ khóa: "consider repeatedly", "final decision".
\item Định vị: Đoạn C: "...engages in more product comparisons... both before and after they make purchasing decisions, and they take longer to decide..." và Đoạn D: "They also tend to brood or ruminate more..."
\item Đối chiếu: "Take longer to decide" (mất nhiều thời gian để quyết định) và "ruminate" (suy ngẫm/nghiền ngẫm) tương ứng với ý nghĩa của việc cân nhắc lặp đi lặp lại (consider repeatedly). Đây là đặc điểm của Maximiser.
\item Đáp án: A
\end{itemize}
\item Câu 31: “participate in the questionnaire of the author”
\begin{itemize}
\item Từ khóa: "participate", "questionnaire", "author".
\item Định vị: Đoạn C: "In particular, we composed a set of statements... Then we had several thousand people rate themselves..."
\item Đối chiếu: Để phân loại ai là Maximiser và ai là Satisficer, các nhà nghiên cứu đã cho hàng nghìn người tham gia đánh giá (rate themselves). Do đó, cả hai nhóm người này đều xuất hiện từ việc tham gia vào bảng câu hỏi nghiên cứu.
\item Đáp án: C
\end{itemize}
\end{itemize}

\cdsection{Phần II: Câu hỏi 32-36 (True/False/Not Given)}

Mục tiêu: Xác định độ chính xác của thông tin dựa trên bài đọc.

\cdstep{1. Chiến lược tiếp cận}
\begin{itemize}
\item Bước 1: Phân tích câu hỏi: Xác định từ khóa (keywords).
\item Bước 2: Đọc dò (Scanning): Tìm đoạn văn chứa thông tin tương ứng.
\item Bước 3: Đối chiếu và Kết luận (True/False/Not Given).
\end{itemize}

\cdstep{2. Phân tích chi tiết từng câu hỏi}
\begin{itemize}
\item Câu 32: “With society’s advancement, more chances make our lives better and happier.”
\begin{itemize}
\item Từ khóa: "society’s advancement", "more chances", "better and happier".
\item Định vị: Đoạn A: "To an extent, the opportunity to choose enhances our lives... Yet recent research strongly suggests that, psychologically, this assumption is wrong... more is not always better than less."
\item Đối chiếu: Bài đọc khẳng định giả định "nhiều lựa chọn hơn là tốt hơn" (more is better) là sai lầm (wrong) về mặt tâm lý. Điều này mâu thuẫn với mệnh đề trong câu hỏi.
\item Đáp án: FALSE
\end{itemize}
\item Câu 33: “There is a difference of findings by different gender classification.”
\begin{itemize}
\item Từ khóa: "difference", "findings", "gender classification".
\item Định vị: Quét toàn bộ bài đọc để tìm các từ liên quan đến giới tính như "gender", "men", "women", "male", "female".
\item Đối chiếu: Bài đọc tập trung vào sự phân loại tâm lý (Maximizers vs Satisficers) chứ không hề đề cập đến sự khác biệt dựa trên giới tính (gender classification).
\item Đáp án: NOT GIVEN
\end{itemize}
\item Câu 34: “The feeling of loss is greater than that of acquisition.”
\begin{itemize}
\item Từ khóa: "feeling of loss", "greater", "acquisition".
\item Định vị: Đoạn F: "Early decision-making research... showed that people respond much more strongly to losses than gains."
\item Đối chiếu: "Respond much more strongly to losses" (phản ứng mạnh hơn với mất mát) đồng nghĩa với "feeling of loss is greater". "Gains" (cái đạt được) đồng nghĩa với "acquisition".
\item Kết luận: Thông tin trùng khớp.
\item Đáp án: TRUE
\end{itemize}
\item Câu 35: “ ‘Good enough’ plays a more significant role in pursuing ‘best’ standards of the maximizer.”
\begin{itemize}
\item Từ khóa: "Good enough", "significant role", "best standards", "maximizer".
\item Định vị: Đoạn G so sánh hai khái niệm: "The latter’s [Satisficer] 'good enough' philosophy... In addition, the 'good enough' standard leads to much less searching... than the maximizer’s 'best' standard."
\item Đối chiếu: Bài đọc tách biệt rõ ràng: "Good enough" là tiêu chuẩn của Satisficer, trong khi "Best" là tiêu chuẩn của Maximizer. Không có thông tin nói rằng "good enough" đóng vai trò trong việc theo đuổi tiêu chuẩn "best" của Maximizer. Thực tế chúng đối lập nhau.
\item Đáp án: FALSE
\end{itemize}
\item Câu 36: “There are certain correlations between the “regret” people and the maximisers.”
\begin{itemize}
\item Từ khóa: "correlations", "regret people", "maximisers".
\item Định vị: Đoạn H: "Not surprisingly, we also found that people with high regret sensitivity tend to be maximizers."
\item Đối chiếu: Cụm từ "tend to be" (có xu hướng là) chỉ ra mối tương quan (correlations) giữa những người nhạy cảm với sự hối tiếc (regret people) và nhóm Maximisers.
\item Kết luận: Thông tin trùng khớp.
\item Đáp án: TRUE
\end{itemize}
\end{itemize}

\cdsection{Phần III: Câu hỏi 37-40 (Multiple Choice)}

Mục tiêu: Chọn đáp án đúng nhất dựa trên nội dung bài đọc.

\cdstep{1. Chiến lược tiếp cận}
\begin{itemize}
\item Bước 1: Đọc câu hỏi để xác định trọng tâm cần tìm.
\item Bước 2: Định vị đoạn văn chứa thông tin liên quan.
\item Bước 3: Loại trừ các phương án sai và chọn phương án khớp ý nhất.
\end{itemize}

\cdstep{2. Phân tích chi tiết từng câu hỏi}
\begin{itemize}
\item Câu 37: “What is the subject of this passage?”
\begin{itemize}
\item Phân tích: Câu hỏi yêu cầu xác định chủ đề chính của toàn bài.
\item \cdblue{\textbf{Đối chiếu}}
\begin{itemize}
\item Đoạn A mở đầu về sự lựa chọn (choose among options).
\item Đoạn B, C, D, E bàn về mối liên hệ giữa việc lựa chọn (Maximizers vs Satisficers) và cảm giác hạnh phúc/không hạnh phúc (unhappy, pleased, well-being).
\end{itemize}
\item Lựa chọn: Đáp án B "choices and Well-being" (Sự lựa chọn và Hạnh phúc/Sự thịnh vượng tinh thần) bao quát chính xác nội dung toàn bài.
\item Đáp án: B
\end{itemize}
\item Câu 38: “According to the conclusion of questionnaires, which of the following statement is correct?”
\begin{itemize}
\item Từ khóa: "conclusion", "questionnaires", "correct".
\item Định vị: Đoạn E thảo luận về kết quả các bảng câu hỏi (questionnaires) về mức độ hạnh phúc (well-being).
\item Đối chiếu: Đoạn E viết: "individuals with high maximization scores experienced less satisfaction with life and were less happy...".
\item Lựa chọn: Đáp án A "maximisers are less happy" phản ánh chính xác kết quả này.
\item Đáp án: A
\end{itemize}
\item Câu 39: “The experimental on theater tickets suggested:”
\begin{itemize}
\item Từ khóa: "experimental", "theater tickets", "suggested".
\item Định vị: Đoạn I mô tả thí nghiệm về vé nhà hát.
\item Đối chiếu: Đoạn I giải thích lý do người mua vé giá trọn gói đi xem nhiều hơn: "...full-price payers would experience more regret if they did not use the tickets because not using the more costly tickets would constitute a bigger loss." (Người trả trọn gói sẽ thấy tiếc hơn nếu không dùng vé vì đó là một mất mát lớn hơn).
\item Lựa chọn: Đáp án D "people will feel regret more when they fail to use a higher-priced purchase" khớp hoàn toàn với ý này.
\item Đáp án: D
\end{itemize}
\item Câu 40: “What is the author’s suggestion on how to increase happiness:”
\begin{itemize}
\item Từ khóa: "author’s suggestion", "increase happiness".
\item Định vị: Câu cuối cùng của đoạn I: "To increase the sense of happiness, we can decide to restrict our options... For example, make a rule to visit no more than two stores..."
\item Đối chiếu: Tác giả gợi ý "restrict our options" (hạn chế các lựa chọn).
\item Lựa chọn: Đáp án C "reduce the choice or option" đồng nghĩa với "restrict our options".
\item Đáp án: C
\end{itemize}
\end{itemize}

\end{document}
